% 06b-supplement.tex 中文版（中英对照用；接在第 5 节之后，同一节内续排）
% 数字由程序从以下文件读出：
%   paper/ablation/results/abl-20261027-s1-engine.json
%   paper/ablation/heldout/results/holdout-20261027-010313.json
\subsection{机制级消融}

第一份补充材料针对第 5.2 节的两个局限：每类三个探针无法达到显著，
而一次改动七个机制则无法归因任何一项。我们从同一批语料出发，
按三种不同的固定置换（SHA-256 \texttt{410085b0\ldots}、\texttt{f58ac2da\ldots}、
\texttt{a690d3ae\ldots}）生成三个 100 轮脚本，每个脚本十六个探针（每类四个），
并让九个臂跑完三个脚本：48 个探针、每类十二个、27 个单元格。
该运行是确定性的、且不使用模型（\texttt{model\_calls=0}），测试装置与流水线在 300/300 轮上一致，
并且探针库在任何臂运行之前就以 SHA-256（\texttt{cdccdde2966e3d49\ldots}）冻结。

\begin{table*}[t]
\centering\footnotesize
\caption{机制级消融（S1）。SFR：按探针类别的设定遗忘率（A 长程、B 被取代、
C 干扰过载、D ACT-R 救援，每类 $n=12$），以及在 48 个探针上的合并值。来源：
\texttt{abl-20261027-s1-engine.json}，字段
\texttt{arms.<arm>.pooled.metrics.m1\_setting\_forgetting.*} 与
\texttt{arms.<arm>.pooled.telemetry.*}；该文件报告 \texttt{valid=true} 与
\texttt{model\_calls=0}。}
\begin{tabular}{@{}l cccc c p{0.42\textwidth}@{}}
\toprule
臂 & A & B & C & D & SFR$_{\text{text}}$ & 观测到的效果与注意点 \\
\midrule
Baseline & 1.00 & 1.00 & 1.00 & 1.00 & 1.00 & 定义性的：一个三轮窗口装不下在第 1 轮埋下的设定 \\
Basic & 0.00 & 0.00 & 0.00 & 1.00 & 0.25 & 只漏掉 ACT-R 类 \\
Full & 0.00 & 0.00 & 0.00 & 0.00 & 0.00 & 参考臂 \\
Full$-$ACT-R & 0.00 & 0.00 & 0.00 & 1.00 & 0.25 & D 恰好回到 Basic；A/B/C 不变 $\Rightarrow$ 该通道是原因 \\
Full$-$anchor & 1.00 & 1.00 & 1.00 & 1.00 & 1.00 & UCA-07 fail-closed：148 个被阻断轮次、\texttt{prompt=None}；这是结构性的，\emph{不是}召回效果 \\
Full$-$gate & 0.00 & 0.00 & 0.00 & 0.00 & 0.00 & 闸门缺席；JI-in 降到 0.00，JI-fp 不变 \\
Full$-$integral & 0.00 & 0.00 & 0.00 & 0.00 & 0.00 & 对保持率无影响；UCA-08 记录从 176 降到 0 \\
Full$-$ushape & 0.00 & 0.00 & 0.00 & 0.00 & 0.00 & 对保持率无影响（决定结果的是成员集合，不是顺序） \\
Full$-$variety & 0.00 & 0.00 & 0.00 & 0.00 & 0.00 & UCA-06 耦合：12 个被阻断轮次；fallback 标记从 12 $\rightarrow$ 0 \\
\bottomrule
\end{tabular}

\vspace{2pt}
\raggedright\scriptsize
每类 $n=12$（三个脚本 $\times$ 四个探针）；对配对结果做精确单侧符号检验，
四类上的 Bonferroni 校正 $\alpha=0.0125$。D 类是一个构造出来的前提条件
（30 次访问，最后一次在 100 天前），因此它的结果关乎前 3 名的成员集合，而不是关乎模型。
Baseline 与 Basic 没有闸门，因此没有可报告的阻断轮数。
\end{table*}

表 3 展示了探针扩展后改变了什么。ACT-R 效应现在可判定：
\textsc{Basic} 忘了全部十二个 D 探针，\textsc{Full} 一个都没忘，
十二个配对结果全部偏向 \textsc{Full}，精确检验在 $\alpha=0.0125$ 下给出
$p=2.4\times10^{-4}$（当时要求至少七个单向翻转）。
\textbf{在只有三个探针时这在结构上不可能——最小可达到的 $p$ 是 $2^{-3}=0.125$——
所以此前的零结果是统计功效的限制，而不是「无效应」的证据。}
去掉 ACT-R 通道会在 D 上精确复现 \textsc{Basic}，且 A/B/C 不受影响，因此归因成立；
但它仍然只支持关于排序的主张，而不支持关于模型行为的主张。

有两个机制臂不可归因，我们照原样报告。从 \textsc{Full} 中去掉锚点，
会使相变监控要求一个永远不会被插入的锚点，于是前置闸触发 UCA-07 并关闭本轮：
148 个被阻断轮次不返回任何 prompt，文本层 SFR 在结构上为 1.00。
这是 fail-closed 的耦合，而不是「锚点影响召回」的证据，并且预注册的 P-S1c 被记为
\emph{失败}，而不是被重新解释。去掉多样性守卫会在十二个冲突轮上触发 UCA-06；
SFR 不变，P-S1d 通过其既定准则，但该臂仍然不能被读作干净的单一机制检验。
去掉闸门保持率不变，只移动越狱通道：JI-in 从 1.00 变为 0.00，
而 JI-fp 仍为 $0/294$，这正是预期的接线确认。去掉积分或去掉 U 型顺序完全不改变 SFR；
只有遥测在动（UCA-08 记录 176$\rightarrow$0）。
我们的结论是：唯一被证明对保持率有效应的机制是 ACT-R 通道；
有两个臂按构造无法与闸门分离；其余机制影响的是审计轨迹而不是选择。

\subsection{留出越狱集}

前面使用的分布内越狱语料是由闸门自己的触发短语构造出来的，因此它 25/25 的拦截率是同义反复。
因此我们改为在两套不同的知识状态下撰写的 40 条改写句上测量闸门。
v1 语料由一位此前读过引擎触发词表的作者撰写（其自述为 semi-blind），并刻意回避原文措辞；
它是对闸门最不利的情形（泄漏上界、拦截下界）。v2 语料由一位从未读过任何词表的作者撰写；
它是盲测的、主要的估计。每套语料有二十条条目，分属五个族（语序倒装、同义替换、跨语言、
指代消解、编码后置）。机器检查确认：两套语料中没有任何一条包含任一词表的字面触发短语。
闸门调用沿用混沌执行器所用的冻结单元级引擎调用，在该调用下只有 UCA-01 能触发；
良性对照保持不变（200 条回放消息加 94 条非越狱脚本轮次）。

\begin{table}[t]
\centering\footnotesize
\caption{留出越狱评测（S2）。JI-held：每套语料 20 条改写句上的 UCA-01 阻断率；
Clopper--Pearson 精确 95\% 区间。JI-fp：294 条良性条目上的误报。来源：
\texttt{holdout-20261027-010313.json}，字段
\texttt{sets.v1|v2.ji\_held.*}、\texttt{ji\_fp.overall.*}、
\texttt{comparison.ji\_held.*}；词表哈希记录在同一文件中。}
\begin{tabular}{@{}l l c c@{}}
\toprule
语料 & 盲测级别 & JI-held [95\% CI] & 被拦截/总数 \\
\midrule
v1（semi-blind） & semi & 0.000 [0.000, 0.168] & 0/20 \\
v2（blind） & blind & 0.000 [0.000, 0.168] & 0/20 \\
差值 v2$-$v1 & --- & 0.000 & 地板效应 \\
JI-fp（良性） & --- & 0.000 [0.000, 0.012] & 0/294 \\
\bottomrule
\end{tabular}
\end{table}

表 4 给出了结果：这个词法闸门对四十条改写句\textbf{一条也没有}拦截。
两套语料里每个族都是 0/4（各族区间 $[0.000, 0.602]$），
因此失败是均匀的，而不是集中在某一种改写风格上。良性对照保持为零，
所以闸门并没有过度拦截；它只是不识别被改写过的意图。
v1--v2 的差值为 0.000，是因为两套语料都处在底板上：当两个条件都被完全漏掉时，
作者知识效应无法估计；我们并不声称「知道词表没有差别」——
只说这一对数据无法测量它。semi-blind 语料仍是上界参照
（有知识的作者可以完全规避），而盲测语料是主要估计，二者在这里一致。

三条限制界定了这个结果。每套语料二十条条目给出的区间很宽，只支持方向性解读；
指代消解族按构造就是单轮的，这使它人为地更难；而且该测量是一次单元级闸门调用，
因为在真实接线中，用户消息从不进入引擎 prompt。
预注册的顺序预测与作者知识预测（P-S2c、P-S2d）在这个底板上不可区分，
因此被记为 null 而不是 pass。词表没有被修改（其 SHA-256 值记录在结果文件中），
所以这个零不是调参的产物。诚实的读法是：前置闸是一份带 fail-closed 接线与可审计性的
词法黑名单，而不是语义防御——在它据以构建的语料上它拦截一切，
而在同一意图的改写句上它什么也拦截不了。
